\documentclass{article}

\usepackage{amsmath} % math stuff
\usepackage{amssymb} % math stuff
\usepackage{array} % equations and stuff
\usepackage{bm} % bold math
%\usepackage{booktabs} % extra table rule options
%\usepackage{caption} % suppressed table numbering; incompatible with revtex, and longtable, I think
\usepackage{comment} % comment environment
%\usepackage{enumitem} % customization of enumeration, itemize, and description
\usepackage[T1]{fontenc} % font encoding for special characters, must also use scalable font package
\usepackage[margin=0.8in]{geometry} % paper sizes and margins (but be careful not to mess up pre-defined pages)
\usepackage{graphicx} % for graphics
%\usepackage{helvet} % default font is the helvetica postscript font
\usepackage[utf8]{inputenc} % special characters in tex input
\usepackage{layouts} % print units like widths
\usepackage{lipsum} % lorem ipsum filler text
\usepackage{lmodern} % scalable font?
\usepackage{longtable} % multi-page tables
\usepackage{makecell} % specify line-breaks in table cells
\usepackage{mathrsfs} % math script font
\usepackage{mhchem} % easier chemical formula
\usepackage{microtype} % allows disabling of ligatures
\usepackage{multicol} % multicolumns
%\usepackage{newcent} % new century schoolbook font
\usepackage{nicefrac}
\usepackage{numprint} % print and format (large) numbers
\usepackage{parskip} % removes paragraph indentation, and adjusts paragraph skip, as well as list items
\usepackage{pdfpages} % add pdf files as pages
%\usepackage{setspace} % adjust text spacing and indents
\usepackage{siunitx} % decimal alignment, and degree symbol
\usepackage{subfigure} % divided figures
%\usepackage{tabu} % extra table options
\usepackage{textcomp} % symbols
\usepackage{threeparttablex} % better footnotes with longtable
\usepackage{titling} % title placement
\usepackage{ulem} % strikethrough text
\usepackage{upgreek} % upright Greek letters
%\usepackage{url} % superceded by hyperref
\usepackage{verbatim} % verbatim environment
\usepackage{xcolor} % colors and color boxes
\usepackage{xspace} % commands that don't eat up white space
\usepackage{hyperref} % links and page setup; should always come last

\hypersetup{
 bookmarks=true,
 colorlinks=true,
 citecolor=blue,
 linkcolor=blue,
 urlcolor=blue,
 pdfstartview={XYZ null null 1.0} % default open view is 100%
}

\DisableLigatures[f,t]{encoding = T1} % disable ff, fi, fl, tt ligatures; without options, it also disables -- = endash
\renewcommand{\arraystretch}{1.0} % extra vertical (and horizontal?) space in tables

% define centered, left- and right-aligned columns with specified widths
\newcommand{\PreserveBackslash}[1]{\let\temp=\\#1\let\\=\temp}
\newcolumntype{C}[1]{>{\PreserveBackslash\centering}p{#1}}
\newcolumntype{L}[1]{>{\PreserveBackslash\raggedright}p{#1}}
\newcolumntype{R}[1]{>{\PreserveBackslash\raggedleft}p{#1}}

\begin{document}

\pagestyle{empty} % don't number pages

% custom title
\begin{center}
{\LARGE Classic Riddler}

\vspace{0.15in}

{\Large 11 March 2022}
\end{center}


\section*{Riddle:}

We're playing a game where you have to pick four whole numbers.
Then I will roll four fair dice.
If any two of the dice add up to any one of the numbers you picked, then you win!
Otherwise, you lose.

For example, suppose you picked the numbers 2, 3, 4 and 12, and the four dice came up 1, 2, 4 and 5.
Then you'd win, because two of the dice (1 and 2) add up to at least one of the numbers you picked (3).

To maximize your chances of winning, which four numbers should you pick?
And what are your chances of winning?


\section*{Solution:}

I enumerate all number choices and dice rolls in the file \texttt{dice\_sums.xlsx}.
I use this to count up which rolls corresponded to successful choices.
I find that the maximally successful choice is 4,6,8,10.
The total probability of winning is $\nicefrac{1,263}{1,296}\approx0.9753$.

By my reasoning, 4,6,8,10 must be the best choice because they are even numbers, and are centered around 7.
7 is the single most likely number to be a sum, with 6,8 and 5,9\dots\ increasing less likely.
So if only a single number had to be picked, 7 would be the best choice.
But for more than one number to pick, it's better to pick numbers that are more spaced out (instead of just 5,6,7,8 for example).
Picking even numbers specifically is good.
With four dice, it's possible to roll all even or all add numbers, and not have an odd sum at all, but it's not possible to avoid an even sum.
So using four distinct even numbers centered around 7 maximizes the probability.
In fact, if six numbers could be picked, 2,4,6,8,10,12 would guarantee a win.

So the solution is \fcolorbox{red}{white}{\textbf{4,6,8,10}} with a winning probability of \fcolorbox{red}{white}{\textbf{\nicefrac{1,263}{1,296}}}\,.

\end{document}